\documentclass[10pt]{article}
\usepackage[utf8]{inputenc}
\usepackage[vietnamese]{babel}
\usepackage[letterpaper, left=155pt, right=45pt, top=40pt, bottom=40pt, headheight=14pt]{geometry}
\usepackage{amsmath,amssymb}
\usepackage{graphicx}
\usepackage{xcolor}
\usepackage[most]{tcolorbox}
\usepackage{fancyhdr}
\usepackage{mathptmx}

% Màu chuẩn giáo trình Stewart
\definecolor{stewartcyan}{RGB}{0, 130, 185}
\definecolor{stewartred}{RGB}{204, 0, 0}

% Thiết lập header và footer
\pagestyle{fancy}
\fancyhf{}
\renewcommand{\headrulewidth}{0pt}
\fancyhead[R]{\textsf{\textbf{MỤC 2.2}\quad Giới hạn của hàm số}\qquad\textbf{\textsf{87}}}

\setlength{\parindent}{0pt}
\setlength{\parskip}{4pt}

\begin{document}

Chẳng hạn, ký hiệu $x \to 5^-$ có nghĩa là chúng ta chỉ xét $x < 5$, và $x \to 5^+$ có nghĩa là chúng ta chỉ xét $x > 5$. Định nghĩa~2 được minh họa trong Hình~6.

\begin{center}
\vspace{2pt}
\noindent\makebox[\linewidth][c]{%
    \llap{\raisebox{4pt}{\makebox[85pt][l]{\textbf{\textsf{HÌNH 6}}}}\hspace{15pt}}%
    \includegraphics[width=\linewidth]{images/p0122_fig1.png}%
}
\vspace{2pt}
\end{center}

Lưu ý rằng Định nghĩa~2 chỉ khác Định nghĩa~1 ở chỗ chúng ta yêu cầu $x$ phải nhỏ hơn (hoặc lớn hơn) $a$. Bằng cách so sánh các định nghĩa này, chúng ta thấy rằng khẳng định sau đây là đúng.

\vspace{3pt}
\begin{tcolorbox}[
    colframe=stewartred,
    colback=white,
    boxrule=0.8pt,
    arc=0mm,
    left=6pt,
    right=6pt,
    top=5pt,
    bottom=5pt
]
\noindent
{\color{stewartred}\setlength{\fboxrule}{0.8pt}\setlength{\fboxsep}{3pt}\fbox{\textbf{\textsf{\color{stewartred}3}}}}%
\hfill
$\displaystyle \lim_{x \to a} f(x) = L \quad\text{khi và chỉ khi}\quad \lim_{x \to a^-} f(x) = L \quad\text{và}\quad \lim_{x \to a^+} f(x) = L$%
\hfill\null
\end{tcolorbox}
\vspace{4pt}

\noindent\textbf{\textsf{\color{stewartcyan}VÍ DỤ 4}}\quad Đồ thị của hàm số $g$ được cho trong Hình~7.

\begin{center}
\vspace{-2pt}
\noindent\makebox[0.60\linewidth][c]{%
    \llap{\raisebox{4pt}{\makebox[85pt][l]{\textbf{\textsf{HÌNH 7}}}}\hspace{\dimexpr 0.20\linewidth + 15pt\relax}}%
    \includegraphics[width=0.60\linewidth]{images/p0122_fig2.png}%
}
\vspace{-2pt}
\end{center}

Dựa vào đồ thị, hãy nêu giá trị (nếu tồn tại) của các giới hạn sau:
\vspace{2pt}

\noindent
\begin{tabular*}{\linewidth}{@{\extracolsep{\fill}}lll}
(a) $\displaystyle \lim_{x \to 2^-} g(x)$ & (b) $\displaystyle \lim_{x \to 2^+} g(x)$ & (c) $\displaystyle \lim_{x \to 2} g(x)$ \\[10pt]
(d) $\displaystyle \lim_{x \to 5^-} g(x)$ & (e) $\displaystyle \lim_{x \to 5^+} g(x)$ & (f) $\displaystyle \lim_{x \to 5} g(x)$
\end{tabular*}

\vspace{6pt}
\noindent\textbf{\textsf{\color{stewartcyan}LỜI GIẢI}}\quad Quan sát đồ thị, ta thấy rằng các giá trị của $g(x)$ tiến dần đến 3 khi $x$ tiến dần đến 2 từ bên trái, nhưng chúng tiến dần đến 1 khi $x$ tiến dần đến 2 từ bên phải. Do đó
\begin{equation*}
\text{(a)}\; \lim_{x \to 2^-} g(x) = 3 \qquad\text{và}\qquad \text{(b)}\; \lim_{x \to 2^+} g(x) = 1
\end{equation*}

(c) Vì giới hạn trái và giới hạn phải khác nhau, nên từ (3) ta kết luận rằng $\displaystyle\lim_{x \to 2} g(x)$ không tồn tại.

\vspace{2pt}
Đồ thị cũng cho thấy rằng
\begin{equation*}
\text{(d)}\; \lim_{x \to 5^-} g(x) = 2 \qquad\text{và}\qquad \text{(e)}\; \lim_{x \to 5^+} g(x) = 2
\end{equation*}

(f) Lần này giới hạn trái và giới hạn phải bằng nhau và vì vậy, theo (3), ta có
\[
\lim_{x \to 5} g(x) = 2
\]

Mặc dù vậy, hãy lưu ý rằng $g(5) \ne 2$.\hfill{\color{stewartcyan}\rule{1.1ex}{1.1ex}}

\end{document}